\chapter{Operational semantics}
\label{ch:opsem}

This section provides a semi-formal definition of the operational semantics of
the language, definining the state and context and how they mutate with every
expression or statement.

\section{Notation and terminology}

\paragraph{Common sets} Some sets will be used throughout the rules
\[
	\begin{split}
		\Expr &= \set{\expr \mid \expr := \text{any expression}} \\*
		\Val  &= \set{\val \mid \val := \text{any value}} \\*
	\end{split}
\]

\paragraph{Expression} The $\expr$ will be used to denote any expression
permitted by the language (e.g. a literal, a binary expression, a function
invocation etc.).\\*
Unless noted otherwise, $\expr \in \Expr$ is implied.

\paragraph{Expression's value} To denote to what value an expression evaluates
to, the $\eval$ symbol will be used, for example: $\expr \eval \val$.\\*
To specify the value along with its type, the notation $\tau(\val)$ will be
used, where $\tau:\Type$, for example: $\expr\eval\Int(\val)$.\\*
Unless noted otherwise, $\val \in \Val$ is implied.

\paragraph{Environments} Two environments will be used:
\begin{itemize}
	\item $\EnvLoc :=$ \textbf{location environment}. Describes the
		mappings/bindings between identifiers or an evaluated expression and
		their "location/address" in memory.
		\[
			\begin{split}
				\EnvLoc: \Id \cup \Expr & \to \mathbb{Z}\\*
				\EnvLoc(x) & \mapsto \addr
			\end{split}
		\]
		An address will usually be denoted by $\addr$
	\item $\EnvStore :=$ \textbf{store environment}. Describes the mappings
		between a location in memory and its data content
		\[
			\begin{split}
				\EnvStore : \mathbb{Z} & \to \Val \\*
				\EnvStore(\addr) & \mapsto \val
			\end{split}
		\]
\end{itemize}

\nb{
	"Location" and "memory" are used loosely here. They do not represent an
	actual computer memory, but rather an abstraction of it.
}

\paragraph{Expression's location} The expression location and its value
mappings in the environment is implied
\[
	\expr\eval\val
	\implies \exists_\addr{\EnvLoc(\expr)=\addr \land \EnvStore(\addr)=\val}
\]

\paragraph{Values identity} In this context, values do not have an identity,
meaning that if two memory locations are mapped to the same value, the value
data is effectively duplicated.

\paragraph{Unmapped values} A slight abuse of notation is used to denote the
lack of mapping in an environment via the $\varnothing$ symbol.\\*
For example: $\EnvLoc(i) = \varnothing, \EnvStore(\addr) = \varnothing$

\paragraph{Environment extension} Like in type inference rules, to denote an
environment extending, modifying or overriding an existing one, this notation
will be used $\EnvLoc\env{i\mapsto\addr}$ where it's defined as follows:
\[
	\EnvLoc\env{i\mapsto\addr}(j) \mapsto \begin{cases}
		\addr      &\text{ if } i = j\\
		\EnvLoc(j) &\text{ if } i \neq j\\
	\end{cases}
\]

\paragraph{The $\alloc()$ function} When a store environment is extended with
a mapping, it's implicit that it is added if it does not exist, and it is
updated if it already exists. In some cases, to improve clarity, the $\alloc()$
function will be used to explicitly state that a new memory "address" is
allocated in the specified environment. The $\alloc()$ function takes a store
environment and returns the newly allocated address.\\*
For example: $\addr = \alloc(\EnvStore)$

% \subparagraph{The $\copi()$ function} The language (almost) never pass values
% by reference, hence many copies are made during an execution. The simplify the
% rules, a $\copi()$ function is used as a shorthand for allocating a new address
% and mapping it a copy of the specified value

\paragraph{Environment modification} If the evaluation of an expression
results in an alteration of an environment, the judgment may be followed by a
"double turnstile" symbol $\Then$ and the new environments.\\*
For example: $\EnvLoc,\EnvStore \then \expr \Then \EnvStore_2$

\clearpage

\subsection{Zero values}
\label{sec:zero_values}

For more info, see~\See{ch:defaults}.

A type's zero value will be described using $\Zero$
\[
	\begin{split}
		\Zero : \Type &\to \Val \\*
		\Zero(\tau)   &\mapsto \val
	\end{split}
\]

Since in some cases the zero value may be overridden, the $\Zero$ function can
be applied on an identifier
\[
	\begin{split}
		\Zero : \Id &\to \Val \\*
		\Zero(i) &\mapsto \begin{cases}
			\val             &\text{ if default overridden } \\*
			\Zero(\EnvSym(i)) &
		\end{cases}
	\end{split}
\]

\subsubsection{Primitive types}

\[
	\begin{split}
		\EnvSym\then \Zero(\Bool)    &= \Bool(\texttt{false}) \\*
		\EnvSym\then \Zero(\Text)    &= \Text(\texttt{""}) \\*
		\EnvSym\then \Zero(\Natural) &= \Natural(\texttt{0}) \\*
		\EnvSym\then \Zero(\Int)     &= \Int(\texttt{0}) \\*
		\EnvSym\then \Zero(\Real)    &= \Real(\texttt{0}) \\*
		\EnvSym\then \Zero(\Type)    &= \Type(\Void) \\*
		\EnvSym\then \Zero(\Void)    &= \Void() \\*
	\end{split}
\]

\subsubsection{Collection types}

\paragraph{List}

\[
	\EnvSym\then \Zero(\List[\nu:=\Type]) = \List(0, \Zero(\nu))
\]

\todo{se eventualmente formalizzo l'operatore \code{**}, aggiornare
dimensione iniziale}

\paragraph{Map}

\[
	\EnvSym\then \Zero(\Map[\kappa:=\Type,\nu:=\Type]) = \Map(\Zero(\kappa), \Zero(\nu))
\]

\subsubsection{Product types}

\todo{riscrivere usando environment per membri (con default override)}

\paragraph{Tuple}

\[
	\infrule{
		\EnvSym\then \forall_{i\in[1,n]}{\tau_i:\Type} \semiq n \geq 0
	}{
		\EnvSym\then \Zero(\tau_1 \times \dotsb \times \tau_n) = \Zero(\tau_1) \times \dotsb \times \Zero(\tau_n)
	}
\]

\paragraph{Record}

\[
	\infrule{
		\EnvSym\then n \geq 0 \semiq \forall_{i\in[1,n]}{\tau_i:\Type, x_i \in \Id}
	}{
		\EnvSym\then \Zero(\tuple{x_1:\tau_1,\dotsc,x_n:\tau_n}) = \tuple{x_1:\Zero(\tau_1), \dotsc, x_n:\Zero(\tau_n)}
	}
\]

\paragraph{Struct}

\text{}

\todo{completare}

\subsubsection{Function}

\[
	\EnvSym\then \Zero(\Function) = ()\to\Void
\]

% \[
% 	\infrule{
% 		\Env\then n \in \mathbb{N} \semiq \forall_{i\in[1,n]}{\tau_i:\Type,x_i\in\Id} \semiq \rho:\Type \\*
% 		\Env\then \phi : (\tau_1,\dotsc,\tau_n)\to\rho
% 	}{
% 		\Env\then \Zero(\phi()) = \Zero(\rho)
% 	}
% \]

\subsubsection{Range}

\[
	\EnvSym\then \Zero(\Range) = \Range(0, 0, 1)
\]

\clearpage

\section{Operational rules}

\subsection{Literals}

\todo{come nelle regole di inferenza?}

\paragraph{Bool}

\[
	\oprule{
		b \in \set{\texttt{true}, \texttt{false}}
	}{
		\EnvLoc,\EnvStore\then \code{$b$} \eval \texttt{Bool($b$)}
	}
\]

\paragraph{Natural}

\[
	\oprule{
		n := \text{ natural literal }
	}{
		\EnvLoc,\EnvStore\then \code{$n$} \eval \texttt{Natural($n$)}
	}
\]

\paragraph{Int}

\[
	\oprule{
		i := \text{ integer literal }
	}{
		\EnvLoc,\EnvStore\then \code{$i$} \eval \texttt{Int($i$)}
	}
\]

\paragraph{Real}

\[
	\oprule{
		r := \text{ real literal }
	}{
		\EnvLoc,\EnvStore\then \code{$r$} \eval \texttt{Real($r$)}
	}
\]


\paragraph{Text}

\[
	\oprule{
		t := \text{ text literal }
	}{
		\EnvLoc,\EnvStore\then \code{$t$} \eval \texttt{Text($t$)}
	}
\]

\todo{variables interpolation?}

\subsection{Identifiers}

\[
	\oprule{
		x \in \Id \semiq \EnvLoc(x) = \addr \semiq \EnvStore(\addr) = \val
	}{
		\EnvLoc,\EnvStore\then \code{$x$} \eval \val
	}
\]

\subsection{Statements}

\paragraph{Redirection (assignment)}

\[
	\oprule[aligned]{
		\EnvLoc,\EnvStore     &\then x \in \Id \semiq \EnvLoc(x) = \addr_x \semiq \EnvStore(\addr_x) = \val_x \\*
		\EnvLoc,\EnvStore     &\then \expr \eval \val \Then \EnvLoc_2,\EnvStore_2 \\*
		\EnvLoc_2,\EnvStore_2 &\then \addr_x' = \alloc(\EnvStore_2) \\*
	}{
		\EnvLoc_2,\EnvStore_2\then \code{$x$ <- $\expr$} \Then \EnvLoc_2\env{x \mapsto \addr_x'},\EnvStore_2\env{\addr_x' \mapsto \val}
	}
\]

\paragraph{Variable declaration}

\[
	\oprule[aligned]{
		\EnvSym           &\then x \in \Id \semiq \tau : \Type \\*
		\EnvLoc,\EnvStore &\then \EnvLoc(x) = \varnothing \\*
		\EnvLoc,\EnvStore &\then \addr = \alloc(\EnvStore)
	}{
		\EnvSym,\EnvLoc,\EnvStore \then \code{$x$ $\tau$}
		\Then \EnvLoc\env{x \mapsto \addr},\EnvStore\env{\addr \mapsto \Zero(\tau)}
	}
\]

\subparagraph{Initialization} A variable can be initialized at declaration
\[
	\oprule[aligned]{
		\EnvSym               &\then x \in \Id \semiq \tau : \Type \\*
		\EnvLoc,\EnvStore     &\then \EnvLoc(x) = \varnothing \\*
		\EnvLoc,\EnvStore     &\then \expr \eval \val \Then \EnvLoc_2,\EnvStore_2 \\*
		\EnvLoc_2,\EnvStore_2 &\then \addr = \alloc(\EnvStore_2) \\*
	}{
		\EnvLoc_2,\EnvStore_2 \then \code{$x$ $\tau$ <- $\expr$}
		\Then \EnvLoc_2\env{x \mapsto \addr},\EnvStore_2\env{\addr \mapsto \val}
	}
\]
It is semantically identical to a variable declaration statement followed by a
different redirection statement.


\subparagraph{Zero override} At (and only at) declaration, the zero value of a
variable can be overridden
\[
	\oprule[aligned]{
		\EnvSym               &\then x \in \Id \semiq \tau : \Type \\*
		\EnvLoc,\EnvStore     &\then \EnvLoc(x) = \varnothing \\*
		\EnvLoc,\EnvStore     &\then \expr \eval \val \Then \EnvLoc_2,\EnvStore_2 \\*
		\EnvLoc_2,\EnvStore_2 &\then \addr = \alloc(\EnvStore_2) \\*
	}{
		\EnvLoc_2,\EnvStore_2 \then \code{$x$ $\tau$ default $\expr$}
		\Then \EnvLoc_2\env{x \mapsto \addr},\EnvStore\env{\addr \mapsto \val}
	}
\]

\subparagraph{Initialization vs Zero override} Initialization and zero
override may seem identical, but they do two different things:
\begin{itemize}
	\item initialization: initializes the variable values with the one provided
		(without overriding its zero value, which remains equal to its type's
		zero value).
	\item zero override: overrides the variable's zero value (to which the
		variable is initialized to).
\end{itemize}
Both can be performed in a single statement
\[
	\oprule[aligned]{
		\EnvSym               &\then x \in \Id \semiq \tau : \Type \\*
		\EnvLoc,\EnvStore     &\then \EnvLoc(x) = \varnothing \\*
		\EnvLoc,\EnvStore     &\then \expr_z \eval \val_z \Then \EnvLoc_2,\EnvStore_2 \\*
		\EnvLoc_2,\EnvStore_2 &\then \expr \eval \val \Then \EnvLoc_3,\EnvStore_3 \\*
		\EnvLoc_3,\EnvStore_3 &\then \addr_z = \alloc(\EnvStore_3) \\*
		\EnvLoc_3,\EnvStore_3 &\then \addr = \alloc(\EnvStore_3) \\*
	}{
		\EnvLoc_3,\EnvStore_3 \then \code{$x$ $\tau$ default $\expr_z$ <- $\expr$}
		\Then \EnvLoc_3\env{x \mapsto \addr},\EnvStore\env{\addr \mapsto \val,\addr_z \mapsto \val_z} \\*
		\hfill\implies \code{$x$} \eval \val \land \Zero(x) = \val_z
	}
\]

\paragraph{Constant declaration}

\[
	\oprule[aligned]{
		\EnvLoc,\EnvStore     &\then x \in \Id \semiq \EnvLoc(x) = \varnothing \\*
		\EnvLoc,\EnvStore     &\then \expr \eval \val \Then \EnvLoc_2,\EnvStore_2 \\*
		\EnvLoc_2,\EnvStore_2 &\then \EnvLoc_2(\expr) = \addr_\expr \semiq \EnvStore_2(\addr_\expr) = \val \\*
	}{
		\EnvLoc_2,\EnvStore_2\then \code{const $x$ <- $\expr$} \Then \EnvLoc_2\env{x \mapsto \addr_\expr}
	}
\]

\todo{rettificare in caso scelgo deep immutability invece che shallow
immutability}

\paragraph{Alias declaration}

\[
	\oprule[aligned]{
		\EnvLoc,\EnvStore     &\then x \in \Id \semiq \EnvLoc(x) = \varnothing \\*
		\EnvLoc,\EnvStore     &\then \expr \eval \val \Then \EnvLoc_2,\EnvStore_2 \\*
		\EnvLoc_2,\EnvStore_2 &\then \EnvLoc_2(\expr) = \addr_\expr \semiq \EnvStore_2(\addr_\expr) = \val \\*
	}{
		\EnvLoc_2,\EnvStore_2\then \code{alias $\expr$ as $x$} \Then \EnvLoc_2\env{x \mapsto \addr_\expr}
	}
\]

\clearpage

\subsection{Patterns}

See \See{ch:patterns} for more details.

\paragraph{Patterns sets} The following sets group patterns to make rules more
readable
\[
	\begin{split}
		\Pattern         &= \PatternTerminal \cup \PatternProduct \cup \PatternIterator\\*
		\PatternTerminal &= \set{\pattern \mid \pattern := \text{any unary pattern or placeholder pattern}}\\*
		\PatternProduct  &= \set{\pattern \mid \pattern := \text{any product pattern}}\\*
		\PatternIterator &= \set{\pattern \mid \pattern := \text{any iterator pattern}}\\*
	\end{split}
\]

\paragraph{The environment set} The $\Env$ will be used to denote a set of
environments.\\*
For example, $\Env_\EnvLoc$ will denote the set of location environments.

\paragraph{Deconstruction function} Given how patterns can be nested, a
recursive function will be used to describe the deconstruction.
\todo{rinominare in Decomposition?}
\[
	\begin{split}
		\decons& : (\Env_\EnvLoc,\Env_\EnvStore,\Val,\Pattern) \to (\Env_\EnvLoc,\Env_\EnvStore) \\*
		\decons&(\EnvLoc,\EnvStore,\val, \pattern) \mapsto \begin{cases}
			(\EnvLoc\env{},\EnvStore\env{})                                                              &\text{ if } \pattern = \code{\_} \\*
			(\EnvLoc\env{a\mapsto\addr},\EnvStore\env{\addr\mapsto\val})                                 &\text{ if } \pattern = \code{$a$} \\*
			\bigcup_{i\in[1,n]}{\decons(\EnvLoc,\EnvStore,\val_i,\pattern_i)}                            &\text{ if } \val:\Product \\*
																										 &\quad\land \pattern = \code{<$\pattern_1$, $\dotsc$, $\pattern_n$>} \\*
			\decons(\EnvLoc,\EnvStore,\val',\pattern')                                                   &\text{ if } \val:\List \\*
																										 &\quad\land \pattern = \code{[$\pattern'$]} \\*
			\decons(\EnvLoc,\EnvStore,\val_k,\pattern_k)\cup\decons(\EnvLoc,\EnvStore,\val_v,\pattern_v) &\text{ if } \val:\Map \\*
																										 &\quad\land \pattern = \code{[$\pattern_k$ => $\pattern_v$]} \\*
		\end{cases}
	\end{split}
\]
In other words, the function tells that:
\begin{itemize}
	\item for a placeholder pattern, the resulting environments will be the
		same.
	\item for a unary pattern, the resulting environments will be extended with
		a new mapping for the specified value.
	\item for product pattern, the resulting environments will be extended with
		the mappings produced from applying the deconstruction function to
		every field of the product value.
	\item for a list pattern, the resulting environments will be extended with
		the mappings produced from applying the deconstruction function to a
		list's entry value.
	\item for a map pattern, the resulting environments will be extended with
		the mappings produced from applying the deconstruction function to a
		map's entry key and value pair.
\end{itemize}

\clearpage

\subsubsection{Terminal patterns}

\paragraph{Placeholder pattern}
\[
	\oprule{
		\EnvLoc,\EnvStore \then \expr \in \Expr \Then \EnvLoc_2,\EnvStore_2
	}{
		\EnvLoc_2,\EnvStore_2 \then \code{$\expr$ -> \_} \Then \EnvLoc_2,\EnvStore_2
	}
\]

\vspace{\baselineskip}

More formally, using the deconstruction function
\[
	\oprule{
		\pattern = \code{\_} \\*
		\EnvLoc,\EnvStore \then \expr\eval\val \Then \EnvLoc_2,\EnvStore_2 \\*
	}{
		\EnvLoc_2,\EnvStore_2 \then \code{$\expr$ -> $\pattern$} \Then \decons(\EnvLoc_2,\EnvStore_2,\val,\pattern) \\*
	}
\]

\paragraph{Unary pattern} The semantic rules are the same of redirections and
variable declarations, because they're secretly the same thing.

\clearpage

\subsubsection{Product patterns}

\paragraph{Positional}

In case only unary patterns are used:
\[
	\oprule[aligned]{
		\EnvLoc,\EnvStore      &\then n \geq 0 \semiq \forall_{i\in[1,n]}{\val_i\in\Val} \semiq \val = \bigtimes_{i \in [1,n]}{\val_i} \\*
		\EnvLoc,\EnvStore      &\then m \leq n \semiq \forall_{j\in[1,m]}{\pattern_j\in\PatternTerminal\setminus\set{\code{\_}}} \semiq \pattern = \code{<$\pattern_1$, $\dotsc$, $\pattern_m$>} \\*
		\EnvLoc,\EnvStore      &\then \expr \eval \val \Then \EnvLoc_2,\EnvStore_2 \\*
		\EnvLoc_2,\EnvStore_2  &\then \forall_{j\in[1,m]}{\addr_j = \alloc(\EnvStore_2)} \\*
	}{
		\EnvLoc_2,\EnvStore_2 \then \code{$\expr$ -> $\pattern$} \Then \EnvLoc_2\env{\forall_{j\in[1,m]}{\pattern_j\mapsto\addr_j}},\EnvStore_2\env{\forall_{j\in[1,m]}{\addr_j\mapsto\val_j}}
	}
\]

\vspace{\baselineskip}

More formally, in case of nested patterns
\[
	\oprule[aligned]{
		\EnvLoc,\EnvStore &\then n \geq 0 \semiq \forall_{i\in[1,n]}{\val_i\in\Val} \semiq \val = \bigtimes_{i \in [1,n]}{\val_i} \\*
		\EnvLoc,\EnvStore &\then m \leq n \semiq \forall_{j\in[1,m]}{\pattern_j\in\PatternTerminal\cup\PatternProduct} \semiq \pattern = \code{<$\pattern_1$, $\dotsc$, $\pattern_m$>} \\*
		\EnvLoc,\EnvStore &\then \expr \eval \val \Then \EnvLoc_2,\EnvStore_2 \\*
	}{
		\EnvLoc_2,\EnvStore_2 \then \code{$\expr$ -> $\pattern$} \Then \decons(\EnvLoc_2,\EnvStore_2,\val,\pattern)
		% \EnvLoc_2\env{\forall_{j\in[1,m]}{\pattern_j\mapsto\addr_j}},\EnvStore_2\env{\forall_{j\in[1,m]}{\addr_j\mapsto\val_j}}
	}
\]

\paragraph{Nominal}

In case only unary patterns are used:
\[
	\oprule[aligned]{
		\EnvLoc,\EnvStore     &\then n \geq 0 \semiq \forall_{i\in[1,n]}{x_i\in\mathcal{N}\subseteq\Id, \val_i\in\Val} \semiq \val = \tuple{x_1:\val_1,\dotsc,x_n:\val_n} \\*
		\EnvLoc,\EnvStore     &\then m \leq n \semiq \forall_{j\in[1,m]}{y_j\in\mathcal{M}\subseteq\mathcal{N},\pattern_j\in\PatternTerminal\setminus\set{\code{\_}}} \\*
		\EnvLoc,\EnvStore     &\then \pattern = \code{<$y_1$:$\pattern_1$, $\dotsc$, $y_m$:$\pattern_m$>} \\*
		\EnvLoc,\EnvStore     &\then \expr \eval \val \Then \EnvLoc_2,\EnvStore_2 \\*
		\EnvLoc_2,\EnvStore_2 &\then \forall_{j\in[1,m]}{\addr_j = \alloc(\EnvStore_2)} \\*
		                      &\EnvLoc_3   = \EnvLoc_2\env{\forall_{j\in[1,m]}{\pattern_j\mapsto\addr_j}}\\*
		                      &\EnvStore_3 = \EnvStore_2\env{\forall_{j\in[1,m]}{\exists_{x'\in\mathcal{N}=y_j\in\mathcal{M}}{\quad\addr_j\mapsto\EnvStore_2(\EnvLoc_2(x'))}}}\\*
	}{
		\EnvLoc_2,\EnvStore_2 \then \code{$\expr$ -> $\pattern$} \Then \EnvLoc_3,\EnvStore_3\\*
	}
\]

\nb{
	A notation abuse is used for $\EnvStore_2(\EnvLoc_2(x'))$, which describes
	the value in memory of the field named $x'$ of the record $\val$.
}

\vspace{\baselineskip}

More formally, in case of nested patterns
\[
	\oprule[aligned]{
		\EnvLoc,\EnvStore  &\then n \geq 0 \semiq \forall_{i\in[1,n]}{x_i\in\mathcal{N}\subseteq\Id, \val_i\in\Val} \semiq \val = \tuple{x_1:\val_1,\dotsc,x_n:\val_n} \\*
		\EnvLoc,\EnvStore  &\then m \leq n \semiq \forall_{j\in[1,m]}{y_j\in\mathcal{M}\subseteq\mathcal{N},\pattern_j\in\PatternTerminal\cup\PatternProduct} \\*
		\EnvLoc,\EnvStore  &\then \pattern = \code{<$y_1$:$\pattern_1$, $\dotsc$, $y_m$:$\pattern_m$>} \\*
		\EnvLoc,\EnvStore  &\then \expr \eval \val \Then \EnvLoc_2,\EnvStore_2 \\*
	}{
		\EnvLoc_2,\EnvStore_2 \then \code{$\expr$ -> $\pattern$} \Then \decons(\EnvLoc_2,\EnvStore_2,\val,\pattern)
	}
\]

\subparagraph{Implicitly named product pattern}

\[
	\oprule[aligned]{
		\EnvLoc,\EnvStore       &\then n \geq 0 \semiq \forall_{i\in[1,n]}{x_i\in\mathcal{N}\subseteq\Id, \val_i\in\Val} \semiq \val = \tuple{x_1:\val_1,\dotsc,x_n:\val_n} \\*
		\EnvLoc,\EnvStore       &\then m \leq n \semiq \forall_{j\in[1,m]}{\pattern_j\in\PatternTerminal\setminus\set{\code{\_}}} \\*
		\EnvLoc,\EnvStore       &\then \pattern = \code{<\_:$\pattern_1$, $\dotsc$, \_:$\pattern_m$>} \\*
		\EnvLoc,\EnvStore       &\then \expr \eval \val \Then \EnvLoc_2,\EnvStore_2 \\*
		\EnvLoc_2,\EnvStore_2   &\then \forall_{j\in[1,m]}{\addr_j = \alloc(\EnvStore_2)} \\*
		                        &\EnvLoc_3   = \EnvLoc_2\env{\forall_{j\in[1,m]}{\pattern_j\mapsto\addr_j}}\\*
		                        &\EnvStore_3 = \EnvStore_2\env{\forall_{j\in[1,m]}{\exists_{x'\in\mathcal{N}=\pattern_j\in\mathcal{M}}{\semiq\addr_j\mapsto\EnvStore_2(\EnvLoc_2(x'))}}}\\*
	}{
		\EnvLoc_2,\EnvStore_2 \then \code{$\expr$ -> $\pattern$} \Then \EnvLoc_3,\EnvStore_3\\*
	}
\]

\subsubsection{Iterator patterns}

\paragraph{List pattern}

\[
	\oprule[aligned]{
		\EnvLoc,\EnvStore &\then n \geq 0 \semiq \expr\eval\glist{\val_1,\dotsc,\val_n} \Then \EnvLoc_2,\EnvStore_2 \\*
						  &\forall_{i\in[1,n]}{(\EnvLoc_i,\EnvStore_i) = \decons(\EnvLoc_2,\EnvStore_2,\val_i,\pattern)}
	}{
		\EnvLoc_2,\EnvStore_2 \then \code{for [$\pattern$] in $\expr$ do} \Then \EnvLoc_i,\EnvStore_i\\*
	}
\]
At each iteration $i$, the $i$-th entry of the list gets deconstructed with
the pattern $\pattern$


\paragraph{Map pattern}

\[
	\oprule[aligned]{
		\EnvLoc,\EnvStore &\then n \geq 0 \semiq \expr\eval\map{\val_{k_1}\mapsto\val_{v_1},\dotsc,\val_{k_n}\mapsto\val_{v_n}} \Then \EnvLoc_2,\EnvStore_2 \\*
						  &\forall_{i\in[1,n]}{(\EnvLoc_i,\EnvStore_i) = \decons(\EnvLoc_2,\EnvStore_2,\val_{k_i},\pattern_k)} \cup \decons(\EnvLoc_2,\EnvStore_2,\val_{v_i},\pattern_v)
	}{
		\EnvLoc_2,\EnvStore_2 \then \code{for [$\pattern_k$ => $\pattern_v$] in $\expr$ do} \Then \EnvLoc_i,\EnvStore_i\\*
	}
\]
At each iteration $i$, the $i$-th key-value pair of the map gets deconstructed
with the patterns $\pattern_k$ and $\pattern_v$ respectively

\clearpage

\subsection{Expressions}

\subsubsection{Operators}

\[
	\oprule[aligned]{
		&\bop = \text{any binary operator} \\*
		\EnvLoc,\EnvStore     &\then \expr_1 \eval \val_1 \Then \EnvLoc_2,\EnvStore_2 \\*
		\EnvLoc_2,\EnvStore_2 &\then \expr_2 \eval \val_2 \Then \EnvLoc_3,\EnvStore_3 \\*
		\EnvLoc_3,\EnvStore_3 &\then \val_3 = \val_1 \bop \val_2 \\*
	}{
		\EnvLoc_3,\EnvStore_3\then \code{$\expr_1$ $\bop$ $\expr_2$} \eval \val_3 \Then \EnvStore_3\env{\alloc(\EnvStore_3) \mapsto \val_3}
	}
\]

\paragraph{Side-effects} Evaluating expressions may cause side effects (e.g.
a function invokation), hence new environments are declared.

\paragraph{Short-circuit evaluation} In some operations (e.g. boolean), some
expressions may be "short-circuited" (they won't be evaluated).\\*
For example: \code{false and someFunction()}, \texttt{someFunction()} will not
be evaluated.


\subsubsection{Invocation}

A function can be invoked using two approaches:
\begin{itemize}
	\item Positional \code{$f$($\expr_1$,$\dotsc$,$\expr_n$)}
	\item Nominal \code{$f$($x_1$:$\expr_1$,$\dotsc$,$x_n$:$\expr_n$)}
\end{itemize}

\paragraph{Positional}

\[
	\oprule[aligned]{
		\EnvSym               &\then n \geq 0 \semiq \forall_{i\in[1,n]}{x_i\in\Id,\tau_i:\Type} \semiq \tau_{n+1}\Type \\*
		\EnvSym               &\then \phi = (x_1:\tau_1,\dotsc,x_n:\tau_n)\to\tau_{n+1} \\*
		\EnvLoc,\EnvStore     &\then f\in\Expr \semiq f \eval \phi \Then \EnvLoc_2,\EnvStore_2 \\*
		\EnvLoc_2,\EnvStore_2 &\then \forall_{i\in[1,n]}{\expr_i\eval\val_i} \Then \EnvLoc_3,\EnvStore_3 \\*
		\EnvLoc_3,\EnvStore_3 &\then \forall_{i\in[1,n]}{\addr_i=\alloc(\EnvStore_3)} \\*
		                      &\EnvLoc_4 = \EnvLoc_3\env{\forall_{i\in[1,n]}{x_i\mapsto\addr_i}} \\*
		                      &\EnvStore_4 = \EnvStore_3\env{\forall_{i\in[1,n]}{\addr_i\mapsto\val_i}} \\*
	}{
		\EnvLoc_3,\EnvStore_3 \then \code{$f$($\expr_1$,$\dotsc$,$\expr_n$)} \eval \val_{\tau_{n+1}} \Then \EnvLoc_4,\EnvStore_4\env{\alloc(\EnvStore_4)\mapsto\val_{\tau_{n+1}}}
	}
\]
First, the expression $f$ to be invoked is evaluated (any expression that
evaluates to a function can be invoked); then all actual parameters $\expr_i$
are evaluated and bound to the respective formal parameter $x_i$.

\paragraph{Nominal}

\[
	\oprule[aligned]{
		\EnvSym               &\then n \geq 0 \semiq \forall_{i\in[1,n]}{x_i\in\mathcal{N}\subseteq\Id,\tau_i:\Type} \semiq \tau_{n+1}:\Type \\*
		\EnvSym               &\then \phi = (x_1:\tau_1,\dotsc,x_n:\tau_n)\to\tau_{n+1} \\*
		\EnvLoc,\EnvStore     &\then f\in\Expr \semiq f \eval \phi \Then \EnvLoc_2,\EnvStore_2 \\*
		\EnvLoc_2,\EnvStore_2 &\then \forall_{i\in[1,n]}{y_i\in\mathcal{N},\expr_i\eval\val_i} \Then \EnvLoc_3,\EnvStore_3 \\*
		\EnvLoc_3,\EnvStore_3 &\then \forall_{i\in[1,n]}{\addr_i=\alloc(\EnvStore_3)} \\*
		                      &\EnvLoc_4 = \EnvLoc_3\env{\forall_{i\in[1,n]}{\exists_{j\in[1,n]}{y_j=x_i \quad x_i\mapsto\addr_i}}} \\*
		                      &\EnvStore_4 = \EnvStore_3\env{\forall_{i\in[1,n]}{\exists_{j\in[1,n]}{y_j=x_i \quad \addr_i\mapsto\val_i}}} \\*
	}{
		\EnvLoc_3,\EnvStore_3 \then \code{$f$($y_1$:$\expr_1$,$\dotsc$,$y_m$:$\expr_m$)} \eval \val_{\tau_{n+1}} \Then \EnvLoc_4,\EnvStore_4\env{\alloc(\EnvStore_4)\mapsto\val_{\tau_{n+1}}}
	}
\]
First, the expression $f$ to be invoked is evaluated (any expression that
evaluates to a function can be invoked); then all actual parameters
$y_i:\expr_i$ are evaluated and bound to the formal parameter whose identifier
$x_i=y_j$, therefore the order of the parameter doesn't matter.

\subparagraph{Invocation labels} A function can declare an "invocation label"
for one or more formal parameters; the invocation label, if defined, has then
to be used instead of the formal parameter name when invoking the function.\\*
The rules are the same for a nominal invocation, except that, if a formal
parameter $x_i$ defines an invocation label $l_i$, then the actual parameter
name $y_j$ is bound to the formal parameter whose label is $l_i=y_j$.

\todo{
	TODO: se $\tau_1 = \Self$ o $\tau_{n+1} = \Self$

	TODO: opzionalità
}


% \[
% 	\oprule[aligned]{
% 		\EnvSym               &\then n \geq 0 \semiq \forall_{i\in[1,n]}{x_i\in\mathcal{N}\subseteq\Id, \tau_i:\Type} \semiq \tau_{n+1}\Type \\*
% 		\EnvSym               &\then \phi = (x_1:\tau_1,\dotsc,x_n:\tau_n)\to\tau_{n+1} \\*
% 		\EnvLoc,\EnvStore     &\then f\in\Expr \semiq f \eval \phi \Then \EnvLoc_2,\EnvStore_2 \\*
% 		\EnvLoc_2,\EnvStore_2 &\then \forall_{i\in[1,n]}{\expr_i\eval\val_i} \Then \EnvLoc_3,\EnvStore_3 \\*
% 		\EnvLoc_3,\EnvStore_3 &\then m \leq n \semiq \forall_{j\in[1,m]}{y_j\in\mathcal{M}\subseteq\mathcal{N}} \\*
% 		\EnvLoc_3,\EnvStore_3 &\then \forall_{i\in[1,n]}{\addr_i=\alloc(\EnvStore_3)} \\*
% 		                      &\EnvLoc_4 = \EnvLoc_3\env{\forall_{y'\in\mathcal{M}}{y'\mapsto\alloc(\EnvStore_3)}} \cup \env{\forall_{x'\in\mathcal{N}\setminus\mathcal{M}}{x'\mapsto\alloc(\EnvStore_3)}} \\*
% 		                      &\EnvStore_4 = \EnvStore_3\env{\forall_{y'\in\mathcal{M}}{\EnvLoc(y')\mapsto\val_i}} \cup \env{\forall_{x'\in\mathcal{N}\setminus\mathcal{M}}{\EnvLoc(x')\mapsto\Zero(x')}}\\*
% 	}{
% 		\EnvLoc_3,\EnvStore_3 \then \code{$f$($y_1$:$\expr_1$,$\dotsc$,$y_m$:$\expr_m$)} \eval \val_{\tau_{n+1}} \Then \EnvLoc_4,\EnvStore_4\env{\alloc(\EnvStore_4)\mapsto\val_{\tau_{n+1}}}
% 	}
% \]


